\documentclass[10pt]{article}
\usepackage[margin=0.9in]{geometry}\usepackage{amsmath,amssymb,booktabs,microtype,xcolor,fancyhdr}\usepackage[hidelinks]{hyperref}
\definecolor{navy}{RGB}{24,55,95}\pagestyle{fancy}\fancyhf{}\lhead{$F_4$ on $V_{26}$}\rhead{Proof and exact verification}\cfoot{\thepage}\setlength{\headheight}{14pt}
\newcommand{\Sg}{\mathcal S}\newcommand{\CT}{\operatorname{CT}}
\title{\color{navy}\textbf{The compact $F_4$ sigma-MGF on $V_{26}$}\\[1mm]\large Proof and exact low-degree matrix verification}
\author{Computational proof companion}\date{17 July 2026}
\begin{document}\maketitle
\begin{abstract}
For the trace-zero Albert module we prove the compact Molien--Weyl expression,
construct the 26-dimensional torus matrices, and exactly verify
$\Sg_{F_4,26}=1+h_2+h_3+O_{\ge4}$.  The one-copy Hilbert series is
$((1-t^2)(1-t^3))^{-1}$.
\end{abstract}

\section{Formula and proof}
The Cauchy identity followed by Haar projection gives
\[
 \boxed{\Sg_{K,V}=\int_K\prod_a\det(I-t_a\rho(g))^{-1}dg
 =\sum_\lambda\dim(\mathbb S_\lambda V)^K s_\lambda(\mathbf t).}
\]
Equivalently, direct determinant expansion proves
$[m_\tau]\Sg_{K,V}=\dim(\bigotimes_i\operatorname{Sym}^{\tau_i}V)^K$.
Weyl integration turns this into
\[
 \frac1{|W|}\CT_z\left[\prod_{\alpha>0}(1-z^\alpha)(1-z^{-\alpha})
 \prod_{a,\mu}(1-t_a z^\mu)^{-m(\mu)}\right].
\]
For $V_{26}$ the weights are the 24 short roots and zero with multiplicity two;
$|W(F_4)|=1152$.  Therefore
\[
 \boxed{\Sg_{F_4,26}=\frac1{1152}\CT_z\left[D_{F_4}(z)
 \prod_a(1-t_a)^{-2}\prod_{\beta\in\Phi_{\rm short}}
 (1-t_a z^\beta)^{-1}\right].}
\]

\section{Invariant tensors}
The trace-zero Albert algebra carries the restricted trace form and the
restricted determinant.  These are respectively a unique symmetric quadratic
and a unique symmetric cubic invariant.  No other Schur type occurs below
degree four, so
\[
 \boxed{\Sg_{F_4,26}=1+h_2+h_3+O_{\ge4}.}
\]
The standard invariant theory of the trace-zero Albert algebra gives
$\mathbb C[V_{26}]^{F_4}=\mathbb C[q,c]$, with generators of degrees two and
three.  Therefore
$\Sg(t,0,\ldots)=((1-t^2)(1-t^3))^{-1}$.  This all-degree invariant-ring
description is an external input; the notebook verifies only the displayed
low-degree terms.

\section{Verification method and evidence}
Starting with the $F_4$ Cartan matrix, the notebook reflects a short simple
root to obtain 24 short roots and appends two zero weights.  It constructs
\[
 M(\theta)=\operatorname{diag}(e^{i\langle\mu,\theta\rangle})_{\mu}\in U(26)
\]
and checks the determinant integrand against its eigenvalue product.  Schur
characters are computed independently by Jacobi--Trudi.  Their Haar integrals
are exact: for a character $f$ invariant under $W$,
\[
 \frac1{|W|}\CT(D_{F_4}f)=\CT\left(f\prod_{\alpha>0}(1-z^{-\alpha})\right),
\]
and the Weyl-denominator identity supplies every coefficient on the right.
The direct monomial-basis check returns
\[
([m_{(2)}],[m_{(1,1)}],[m_{(3)}],[m_{(2,1)}],[m_{(1,1,1)}])
=(1,1,1,1,1).
\]

\begin{center}\small
\begin{tabular}{@{}c|rrrrrrr@{}}\toprule
$\lambda$&$\varnothing$&$(1)$&$(2)$&$(1,1)$&$(3)$&$(2,1)$&$(1,1,1)$\\\midrule
proposed&1&0&1&0&1&0&0\\
exact Weyl check&1&0&1&0&1&0&0\\\bottomrule
\end{tabular}
\end{center}
The generated root count is 48, the matrix dimension is 26, the unitarity
residual is $2.8\times10^{-16}$, and the two integrand evaluations agree to
$1.3\times10^{-14}$.  Every degree-at-most-three coefficient passes.

\paragraph{Reproducibility.} Run \texttt{marimo\_notebooks/lie\_group\_f4.py} with marimo.
\section*{External invariant-theory source}
See T.~A.~Springer and F.~D.~Veldkamp, \emph{Octonions, Jordan Algebras and
Exceptional Groups}, Springer Monographs in Mathematics (2000), for the
Albert-algebra realization and its norm invariants.
\end{document}
